\subsection{半单模中的标量扩张}

命题 8. —— 设 A 是一个 K-代数，M 是一个 A-模，L 是域 K 的一个扩张。用 D 表示 M 的交换子，用 Z 表示 D 的中心。

\begin{enumerate}
\def\labelenumi{\alph{enumi})}
\item
  假设 \(A_{(L)}\) -模 \(M_{(L)}\) 是单的（相应地，同型的、半单的）。则 A-模 M 是单的（相应地，同型的、半单的）。
\item
  假设 A-模 M 是半单的，且 M 或 L 在 K 上是有限维的。\(\mathrm{A}_{(\mathrm{L})}\) -模 \(\mathrm{M}_{(\mathrm{L})}\) 是半单的当且仅当环 \(\mathrm{Z}_{(\mathrm{L})}\) 是约化的。\(\mathrm{A}_{(\mathrm{L})}\) -模 \(\mathrm{M}_{(\mathrm{L})}\) 是同型且非零的当且仅当环 \(\mathrm{Z}_{(\mathrm{L})}\) 是一个整环。
\item
  假设 A-模 M 是单的。\(\mathrm{A}_{(\mathrm{L})}\) -模 \(\mathrm{M}_{(\mathrm{L})}\) 是半单的（相应地，同型的、单的）当且仅当环 \(\mathrm{D}_{(\mathrm{L})}\) 是半单的（相应地，单的、一个体）。
\end{enumerate}

断言 a) 是命题 1（VIII, 第 211 页）的特殊情形，断言 b) 是命题 5（VIII, 第 219 页）的特殊情形，而断言 c) 是 VIII, 第 215 页的推论 1 的特殊情形。

推论 1. —— a) 假设 A-模 M 是半单的，K 的扩张 L 是可分的，且 M 或 L 在 K 上是有限维的。则 \(\mathrm{A}_{(\mathrm{L})}\) -模 \(\mathrm{M}_{(\mathrm{L})}\) 是半单的。

\begin{enumerate}
\def\labelenumi{\alph{enumi})}
\setcounter{enumi}{1}
\tightlist
\item
  假设 A-模 M 是单的且其交换子等于 K。则 \(\mathrm{A}_{(\mathrm{L})}\) -模 \(\mathrm{M}_{(\mathrm{L})}\) 是单的。
\end{enumerate}

断言 a) 由 VIII, 第 221 页的推论得出。断言 b) 是命题 8, c) 的特殊情形。

推论 2. —— 设 L 是域 K 的一个扩张。用 Z 表示 K-代数 A 的中心。

\begin{enumerate}
\def\labelenumi{\alph{enumi})}
\item
  若 L-代数 \(A_{(L)}\) 是半单的，则 K-代数 A 是半单的。
\item
  假设 K-代数 A 是半单的，且 L 或 A 在 K 上是有限维的。L-代数 \(\mathrm{A}_{(\mathrm{L})}\) 是半单的当且仅当环 \(\mathrm{Z}_{(\mathrm{L})}\) 是约化的；特别地，当 L 是 K 的一个可分扩张时即是这种情形。L-代数 \(\mathrm{A}_{(\mathrm{L})}\) 是单的当且仅当环 \(\mathrm{Z}_{(\mathrm{L})}\) 是一个整环；特别地，当 A 的中心等于 K 时即是这种情形。
\end{enumerate}

断言 a) 与 b) 由应用于 A-模 \(\mathrm{A}_s\) 的命题 8, a) 与 b) 得出。

命题 9. —— 设 A 是一个 K-代数，L 是 K 的一个可分扩张。

\begin{enumerate}
\def\labelenumi{\alph{enumi})}
\item
  若 M 是一个无根的 A-模，则 \(\mathrm{A}_{(\mathrm{L})}\) -模 \(\mathrm{M}_{(\mathrm{L})}\) 是无根的。
\item
  若 K-代数 A 是无根的，则 L-代数 \(\mathrm{A}_{(\mathrm{L})}\) 是无根的。
\end{enumerate}

我们来证明断言 a)。设 M 是一个无根的 A-模。我们把 M 与它在 \(\mathrm{M}_{(\mathrm{L})}\) 中的典范像等同起来。设 N 是 M 的一个极大子模。由于 A-模 M/N 是单的，由 VIII, 第 218 页的命题 4 得出，\(\mathrm{A}_{(\mathrm{L})}\) -模 \((\mathrm{M}/\mathrm{N})_{(\mathrm{L})} = \mathrm{M}_{(\mathrm{L})}/\mathrm{N}_{(\mathrm{L})}\) 是无根的，因此由 VIII, 第 152 页的推论 1, c) 有 \(\mathfrak{R}_{\mathrm{A}_{(\mathrm{L})}}(\mathrm{M}_{(\mathrm{L})}) \subset \mathrm{N}_{(\mathrm{L})}\)。另一方面，由 II, §7, No.~7, 第 306 页的命题 14 的推论得出，当 N 遍历 M 的极大子模所成之集时，诸 \(\mathrm{N}_{(\mathrm{L})}\) 的交化为 0。因此，\(\mathrm{A}_{(\mathrm{L})}\) -模 \(\mathrm{M}_{(\mathrm{L})}\) 是无根的。

断言 b) 由应用于 A-模 \(\mathrm{A}_s\) 的断言 a) 得出。

命题 10. —— 设 A 是一个 K-代数，L 是 K 的一个扩张。设 M 是一个 A-模。

\begin{enumerate}
\def\labelenumi{\alph{enumi})}
\tightlist
\item
  我们作以下两个假设之一：
\end{enumerate}

\begin{enumerate}
\def\labelenumi{(\roman{enumi})}
\item
  A-模 M 是有限生成的，且 L 在 K 上是代数的。
\item
  环 A 是左阿廷的。
\end{enumerate}

则我们有包含关系

\[
\mathfrak {R} _ {\mathrm{A}} (\mathrm{M}) _ {(\mathrm{L})} \subset \mathfrak {R} _ {\mathrm{A} _ {(\mathrm{L})}} (\mathrm{M} _ {(\mathrm{L})}).
\]

\begin{enumerate}
\def\labelenumi{\alph{enumi})}
\setcounter{enumi}{1}
\tightlist
\item
  若 L 是 K 的一个可分扩张，则我们有包含关系
\end{enumerate}

\[
\mathfrak {R} _ {\mathrm{A} _ {(\mathrm{L})}} (\mathrm{M} _ {(\mathrm{L})}) \subset \mathfrak {R} _ {\mathrm{A}} (\mathrm{M}) _ {(\mathrm{L})}.
\]

我们来证明断言 a)。考虑情形 (i)。首先假设 L 在 K 上次数有限。则 A-模 \(M_{(L)}\) 是有限生成的。用 f 表示从 A 到 \(A_{(L)}\) 的典范同态；环 \(A_{(L)}\) 由其中心与 \(f(A)\) 的并生成。因此我们可以把 VIII, 第 174 页的命题 3 应用于 \(A_{(L)}\) -模 \(M_{(L)}\)。这给出包含关系 \(\mathfrak{R}_{\mathrm{A}}(\mathrm{M}_{(\mathrm{L})}) \subset \mathfrak{R}_{\mathrm{A}_{(\mathrm{L})}}(\mathrm{M}_{(\mathrm{L})})\)，更有 \(\mathfrak{R}_{\mathrm{A}}(\mathrm{M})_{(\mathrm{L})} \subset \mathfrak{R}_{\mathrm{A}_{(\mathrm{L})}}(\mathrm{M}_{(\mathrm{L})})\)（VIII, 第 152 页，推论 1）。

现在来处理一般情形。设 \(x_{1},\ldots,x_{n}\) 是生成 A-模 M 的元素，x 是 \(\mathfrak{R}_{\mathrm{A}}(\mathrm{M})\) 的一个元素。设 \(a_{1},\ldots,a_{n}\) 是 \(\mathrm{A}_{(\mathrm{L})}\) 的元素；由于 L 在 K 上是代数的，存在 K 的一个包含在 L 中的有限扩张 \(L'\)，使得诸 \(a_{i}\) 属于 \(\mathrm{A}_{(\mathrm{L}')}\)。由前所述，x 属于 \(\mathrm{A}_{(\mathrm{L}')}\) -模 \(\mathrm{M}_{(\mathrm{L}')}\) 的根；由 VIII, 第 153 页的推论得出，诸元素 \(x_{i}+a_{i}x\)（\(1\leqslant i\leqslant n\)）生成 \(\mathrm{A}_{(\mathrm{L}')}\) -模 \(\mathrm{M}_{(\mathrm{L}')}\)，从而生成 \(\mathrm{A}_{(\mathrm{L})}\) -模 \(\mathrm{M}_{(\mathrm{L})}\)。由同一推论，x 属于 \(\mathrm{A}_{(\mathrm{L})}\) -模 \(\mathrm{M}_{(\mathrm{L})}\) 的根。

现在考虑情形 (ii)。设 \(\mathfrak{r}\) 是 A 的根，于是 A-模 M 的根等于 \(\mathfrak{rM}\)（VIII, 第 174 页，推论）。A 的根 \(\mathfrak{r}\) 是 A 的一个幂零双边理想（VIII, 第 173 页，命题 1）；因此 \(\mathfrak{r}_{(L)}\) 是 \(A_{(L)}\) 的一个幂零双边理想。由此得出 \(\mathfrak{r}_{(L)} \subset \Re(A_{(L)})\)（VIII, 第 156 页，定理 1），且 VIII, 第 158 页的命题 6 蕴含 \(\Re(A_{(L)})M_{(L)} \subset \Re_{A_{(L)}}(M_{(L)})\)。于是我们有 \(\Re_A(M) = \mathfrak{rM} \subset \Re_{A_{(L)}}(M_{(L)})\)，这就完成了断言 a) 的证明。

模 \(\mathrm{M}/\mathfrak{R}_{\mathrm{A}}(\mathrm{M})\) 是无根的。若 L 是 K 的一个可分扩张，则由 VIII, 第 223 页的命题 9 得出，\(\mathrm{A}_{(\mathrm{L})}\) -模 \((\mathrm{M}/\mathfrak{R}_{\mathrm{A}}(\mathrm{M}))_{(\mathrm{L})}\) 是无根的。因此，我们有包含关系

\[
\mathfrak {R} _ {\mathrm{A} _ {(\mathrm{L})}} (\mathrm{M} _ {(\mathrm{L})}) \subset \mathfrak {R} _ {\mathrm{A}} (\mathrm{M}) _ {(\mathrm{L})}.
\]

推论. —— 设 L 是 K 的一个可分扩张。若 L 在 K 上是代数的，或者环 A 是左阿廷的，则我们有 \(\mathfrak{R}(A_{(L)}) = \mathfrak{R}(A)_{(L)}\)。

这是命题 10 中 \(\mathrm{M} = \mathrm{A}_s\) 的情形。

